% ==========================================================================
% DRAFT — Option B reframing (EIC W1, perspective W1–W5; panel 2026-11-07)
% Replacement passages for the ICLR build of main.tex. Not \input anywhere.
%
% Goal: lead with the result a practitioner or auditor can use (sampler knobs
% reweight training examples, they do not protect them; late-time field error,
% not loss, is what moves a learned sampler off the data), and present the
% composition result as a boundary marker of the positive theorems.
%
% Review items this draft is meant to close or advance:
%   R1  rescope practical framing of the composition result  -> abstract, intro, Implications
%   R2  Thm 5.1 takeaway restricted to identifying labels    -> kept, wording unchanged in substance
%   R6  inductive-bias positioning                           -> intro sentence kept (lines 181–189)
%   R8  "collapse" vs endpoint atom                          -> abstract keeps "approach the endpoint"
%   R9  CIFAR pairs as a consistency check                   -> abstract keeps that wording
%   perspective W1 "new samples" = blends                    -> intro list + practice paragraph
%   perspective W2 Ethics not usable                         -> Ethics rewrite
%   perspective W3 fixed-N                                   -> practice paragraph, last item
%   perspective W4 practitioner knobs                        -> abstract no longer opens on them
%   perspective W5 endpoint noise is a smoothed bootstrap    -> practice paragraph, first item
%
% TODO before merging:
%   [T1] refs.bib has no differential-privacy reference. Add one (e.g. Dwork & Roth 2014,
%        "The Algorithmic Foundations of Differential Privacy") or drop the DP clause.
%   [T2] "~10^5" quotes K = (2.1 ± 2.3)×10^5 at 1-t = 1e-3 (App. measured-transfer, EXP-1, h = 0).
%        The sd is as large as the mean; keep "orders of magnitude" if you prefer not to quote it.
%   [T3] The N^{-2/d} remark is copied from the non-ICLR branch after Prop. atomicity
%        (main.tex lines 394–397); either enable that branch in the ICLR build or keep it here.
%   [T4] Check the abstract length against the ICLR limit after editing.
%   [T5] If S10 gives a clean per-time profile, add one clause to the third practice item
%        pointing to App. measured-transfer for the CIFAR measurement.
% ==========================================================================


% --------------------------------------------------------------------------
% 1. ABSTRACT  (replaces main.tex lines 101–122)
% --------------------------------------------------------------------------
\begin{abstract}
Label noise, coarse label resolution and classifier-free guidance modify the field a conditional
flow-matching sampler integrates, and some of them are used to reduce the replication of training
data. Applied to the exact empirical field at positive label bandwidth, none of them does: they
decide which training example is returned, not whether one is. Label noise reweights the training
examples with weights computable from the training set; positive mixtures of exact fields, which
include averaging over labels and negative guidance, and guidance with a fixed nonnegative strength
or a bounded schedule of bounded variation end on a training example from almost every initial
point, and every trajectory converges. What moves a learned sampler off the training set is
therefore its field error late in sampling: at an identifying label a learned endpoint lies within
the limiting scaled field error of its training example, a quantity the regression loss does not
control, and in our runs the error met along trajectories exceeds the error on the training law by
orders of magnitude near $t=1$. At positive bandwidth we prove only a local version. The positive
theorems stop at one composition: guidance applied to a label-averaged field reaches a blend of
training examples with positive probability, for two examples whose averaged priors have log-odds
$\pm2B$ exactly when the guidance strength exceeds $1/\sinh^2B$, stably under perturbation; a signed
posterior covariance decides the general case. We present it as a boundary marker, not a practical
mechanism: we know of no trained model that averages its field over labels, the basin sits in a
layer of width $(1-t)^2$ at the tie, and standard velocity networks do not reproduce it, while
networks given a coordinate that resolves this scale at the known tie do. Exact computations agree
with the threshold, also in a consistency check on CIFAR-10 image pairs. At a finite cutoff,
trained networks from small MLPs to a 35.75M-parameter U-Net approach the hard-conditioning
endpoint as predicted, although at image scale some samples stay off the training set and the
variance departs from the label-smoothing reference.
\end{abstract}


% --------------------------------------------------------------------------
% 2. INTRODUCTION, end of the second paragraph and contribution list
%    (replaces main.tex lines 144–177, up to "\ificlr\else{} Non-training endpoints")
% --------------------------------------------------------------------------
varies in time~\citep{kynkaanniemi2024interval} and is sometimes withheld or reversed early in
sampling to reduce memorisation~\citep{jain2025attraction}. We ask which of these
modifications can, by themselves, move the endpoints of the exact field off the training set,
and what a learned field must do to inherit the answer.
\begin{itemize}\itemsep1pt
\item \emph{Sampler-side modifications reweight the training examples}
  (Section~\ref{sec:intervention-atomicity}). Label noise returns training examples with weights
  $p^{(h)}(y)$ computable from the training set (Theorem~\ref{thm:endpoint}). Positive mixtures of
  prior-indexed exact fields (\emph{positive field mixtures}) with strictly positive priors, and
  nonnegative guidance with a fixed strength or a bounded schedule of bounded variation, end on a
  training example for almost every initial point (Theorems~\ref{thm:posmix}
  and~\ref{thm:schedule}). The exceptional set is null but can be a hyperplane
  (Proposition~\ref{prop:hyperplane}).
\item \emph{Learned fields leave only through late-time error} (Section~\ref{sec:transfer}). A
  learned endpoint is within $\limsup_{t\uparrow1}(1-t)|\Delta|$ of its training example, where
  $\Delta$ is the field error along the trajectory (Theorem~\ref{thm:transfer}). Small regression
  loss does not control this quantity: on EXP-1 the error along trajectories
  exceeds the training-law error by a factor that grows roughly like $(1-t)^{-2}$
  (Appendix~\ref{app:measured-transfer}). A field with integrable Lipschitz constant up to $t=1$
  has an endpoint law without atoms (Proposition~\ref{prop:regular}).
\item \emph{One composition marks where these results stop} (Section~\ref{sec:composition}).
  Guidance applied to a label-averaged field can converge to a circumcentre of training examples,
  a blend of them, with positive probability; for two examples whose averaged priors have log-odds
  $\pm2B$, the threshold is exactly $w>1/\sinh^2B$, and the effect is stable under perturbation of
  the priors, mixture weights and guidance strength (Theorem~\ref{thm:composition},
  Proposition~\ref{prop:composition-open}). For hyperbolic profiles in any dimension, the
  eigenvalue signs of a signed posterior covariance matrix distinguish open attracting basins
  from null basins; a complete classification also needs a condition on the faces of the tied
  atoms (Theorem~\ref{thm:spectral}). Averaging and guidance commute, so the composition, not the
  order, decides the support. We read this as a boundary marker, not as a mechanism of
  standard models: the basin lives at the scale $(1-t)^2$ at the tie, only networks given an
  input that resolves that scale reproduce it (Section~\ref{sec:intervention-tests}), and a label
  bottleneck averages priors rather than fields (Proposition~\ref{prop:bottleneck}).
\end{itemize}
No single intervention analysed here (Table~\ref{tab:interventions}) moves the endpoints of the
field that training defines off the training examples, and the one composition that does reaches
blends of them, through a boundary layer that standard networks do not resolve.

\paragraph{What this means in practice.} Four consequences follow without new assumptions.
\emph{(i) The knobs are not privacy mitigations.} At the exact optimum, label noise, label
averaging, negative guidance and bounded guidance schedules only reweight training examples, and
the blends reached by the composition are convex combinations of training images, which an audit
should count as leakage rather than as new samples. Changing the support needs a mechanism such as
endpoint noise (Proposition~\ref{prop:tgtnoise}), whose endpoint law is a Gaussian kernel density
estimate around the weighted training examples, so sampling from it is a smoothed bootstrap; it carries no formal privacy guarantee unless its
randomisation is calibrated in the sense of differential privacy~\citep{TODO-dwork2014}.
\emph{(ii) Which record is returned is predictable.} Under label noise the returned record and
its probability, $p_i^{(h)}(y)$, are computable from the training set and the query label
(Theorem~\ref{thm:endpoint}), so an auditor can rank, for any label, the records an exact-optimum
model returns; under guidance the support is the same but the weights have no closed form.
\emph{(iii) Monitor late-time field error, not the loss.} Because the exact field is explicit for
a fixed training set, $(1-t)|\Delta|$ can be evaluated along sampling trajectories near $t=1$;
by Theorem~\ref{thm:transfer} this, and not the regression loss, is what bounds the distance of an
endpoint from its training example at an identifying label.
\emph{(iv) The statements concern support at fixed $N$.} The floor in~\eqref{eq:w2floor} is the
mean squared distance from the posterior to the nearest training point, which typically decreases
like $N^{-2/d}$, so atomic support is compatible with a small distance to the posterior at large
$N$ and says nothing on its own about how well the weights match it.


% --------------------------------------------------------------------------
% 3. CONCLUSION, Implications (replaces main.tex lines 1135–1147, ICLR branch)
% --------------------------------------------------------------------------
\emph{Implications.} For a practitioner, label noise, label averaging and guidance decide which
training examples the exact field returns, not whether it returns one, so at the exact optimum
none of them protects training data; under label noise the returned records are even predictable
from the training set. At an identifying label, what moves a trained sampler away from the
training examples is the error of the learned field near $t=1$ (Theorem~\ref{thm:transfer}),
which the regression loss does not control, so late-time field error is the quantity to measure.
At positive bandwidth, as for our off-atom image samples, this attribution rests only on ruling
out the sampler. The composition result is a boundary marker of these statements, not a mechanism
we expect in standard models: its basin is set at the scale $(1-t)^2$, which standard velocity
networks did not resolve in our runs (Proposition~\ref{prop:regular}), and a label bottleneck does
not produce the averaged field (Proposition~\ref{prop:bottleneck}). Open are nonhyperbolic
profiles, many atoms near a candidate, learned composition beyond $d\le2$, and the off-atom
samples of trained image samplers.


% --------------------------------------------------------------------------
% 4. ETHICS STATEMENT (replaces main.tex lines 1173–1180)
% --------------------------------------------------------------------------
\subsection*{Ethics statement}
This work studies memorisation in conditional flow-matching models: when the exact empirical field
and its common modifications return training examples. The consequence for sensitive data is
negative: at the exact optimum, label noise, label averaging, negative guidance and bounded
guidance schedules are not privacy mitigations, because they only reweight the training examples,
and the blends reached by the one composition that leaves them are combinations of training
records. Under label noise the returned record is also predictable from the training set and the
query label, which helps auditors and, equally, anyone attempting extraction. Moving the support
requires a change such as endpoint noise, and a formal guarantee requires randomisation calibrated
in the sense of differential privacy, which we do not analyse. All experiments use synthetic data
and the public CIFAR-10 and MNIST datasets; no personal or human-subject data are involved.
